% Part: many-valued-logic
% Chapter: three-valued-logics
% Section: introduction

\documentclass[../../../include/open-logic-section]{subfiles}

\begin{document}

\olfileid{mvl}{thr}{int}

\olsection{引言}

如果在真与假之外仅仅再增加一个真值~$\Undef$，我们就得到了一种三值逻辑。
虽然只是增加了一个真值，但为 $\lnot$、$\land$、$\lor$ 和 $\lif$ 定义真值函数的可能方式却相当多。
这样一来，一种逻辑可以采用这些真值函数的任意组合，此外还可以选择仅将 $\True$ 作为指定值，或者将 $\True$ 和~$\Undef$ 均作为指定值。

在此，我们介绍若干最著名的三值逻辑，包括它们的动机及部分性质。

\end{document}